\ifdefined\gkver\else\def\gkver{student}\fi
\documentclass[\gkver]{gaokaozhenti}
\begin{document}

\chapter{2025年天津}

\begin{enumerate}
\item
%% number: 1
%% paperName: 2025年天津高考第 1 题 5 分
%% typeId: 1
%% type: 单选题
%% chapter: 曲线运动
%% point: 人造地球卫星
%% method: 无
%% score: 5
%% degree: 750
%% duplicateId: 0
%% body:
2025 年 5 月我国成功发射通信技术试验卫星十九号，若该系列试验卫星中 A、B 两颗卫星均可视为绕地球做匀速圆周运动，轨道半径 $r_{A} > r_{B}$，则卫星 A 比 B\xzanswer{A}

\fourchoices[answer=A]
{线速度小、角速度小}
{线速度小、运行周期小}
{加速度大、角速度大}
{加速度大、运行周期大}

%% answer: A
%% memo:
\memoanswer{
根据题意，由万有引力提供向心力有
$\frac{GMm}{r^{2}} = m\frac{v^{2}}{r} = m\omega^{2}r = m\frac{4\pi^{2}}{T^{2}}r = ma$

解得 $v = \sqrt{\frac{GM}{r}}$，$\omega = \sqrt{\frac{GM}{r^{3}}}$，$T = \sqrt{\frac{4\pi^{2}r^{3}}{GM}}$，$a = \frac{GM}{r^{2}}$

由于轨道半径 $r_{A} > r_{B}$

可得 $v_{A} < v_{B}$，$\omega_{A} < \omega_{B}$，$T_{A} > T_{B}$，$a_{A} < a_{B}$

故选 A。
}

\item
%% number: 2
%% paperName: 2025年天津高考第 2 题 5 分
%% typeId: 1
%% type: 单选题
%% chapter: 原子和原子核
%% point: 天然放射性
%% method: 无
%% score: 5
%% degree: 750
%% duplicateId: 0
%% body:
许多放射性元素要经过多次衰变才能达到稳定，衰变过程中既有 $\alpha$ 衰变也有 $\beta$ 衰变。下列说法正确的是\xzanswer{B}

\fourchoices[answer=B]
{元素发生 $\alpha$ 衰变后质量数增加}
{衰变过程中存在质量亏损}
{低温会增大放射性元素的半衰期}
{$\beta$ 衰变说明原子核内存在电子}

%% answer: B
%% memo:
\memoanswer{
A．$\alpha$ 衰变会释放出氦核 \ce{^{4}_{2}He}（质量数 4），母核质量数减少 4，故 A 错误；

B．衰变释放能量，根据质能方程，总质量减少（质量亏损），故 B 正确；

C．半衰期由核内部结构决定，与温度无关，故 C 错误；

D．$\beta$ 衰变的电子是中子转化为质子时产生的，并非核内原有电子，故 D 错误。

故选 B。
}

\item
%% number: 3
%% paperName: 2025年天津高考第 3 题 5 分
%% typeId: 1
%% type: 单选题
%% chapter: 曲线运动
%% point: 圆周运动的应用
%% method: 无
%% score: 5
%% degree: 650
%% duplicateId: 0
%% body:
一种名为“飞椅”的游乐设施如图所示，该设施中钢绳一端系着座椅，另一端系在悬臂边缘。绕竖直轴转动的悬臂带动座椅在水平面内做匀速圆周运动，座椅可视为质点，则某座椅运动一周的过程中\xzanswer{B}

\onepicture[w=5.00cm, h=4.51cm]{figs/03.jpg}

\fourchoices[answer=B]
{动量保持不变}
{所受合外力做功为零}
{所受重力的冲量为零}
{始终处于受力平衡状态}

%% answer: B
%% memo:
\memoanswer{
座椅在水平面内做匀速圆周运动，速度大小不变，方向改变

A．根据 $p = mv$ 可知动量大小不变，方向改变，故 A 错误；

B．速度大小不变，则座椅的动能不变，根据动能定理可知所受合外力做功为零，故 B 正确；

C．根据 $I_{G} = mgt$ 可知所受重力的冲量不为零，故 C 错误；

D．座椅在水平面内做匀速圆周运动，一定有向心加速度，所以不是处于受力平衡状态，故 D 错误。

故选 B。
}

\item
%% number: 4
%% paperName: 2025年天津高考第 4 题 5 分
%% typeId: 1
%% type: 单选题
%% chapter: 气体性质
%% point: 热力学第一定律
%% method: 无
%% score: 5
%% degree: 650
%% duplicateId: 0
%% body:
如图所示，一定质量的理想气体可经三个不同的过程从状态 A 变化到状态 C，则\xzanswer{C}

\onepicture[w=4.50cm, h=3.97cm]{figs/04.png}

\fourchoices[answer=C]
{AC 和 ADC 过程，外界对气体做功相同}
{ABC 和 ADC 过程，气体放出的热量相同}
{在状态 A 时和在状态 C 时，气体的内能相同}
{在状态 B 时和在状态 D 时，气体分子热运动的平均动能相同}

%% answer: C
%% memo:
\memoanswer{
A．根据 $W = pV$，AC 过程的压强总比 ADC 过程的压强大，则 AC 过程外界对气体做功多，故 A 错误；

B．同理可知 ABC 过程外界对气体做功比 ADC 过程大，根据热力学第一定律 $\Delta U = W + Q$ 可知 ABC 过程气体放出的热量多，故 B 错误；

C．根据理想气体状态方程结合题图可知状态 A 和状态 C 温度相等，则气体内能相同，故 C 正确；

D．根据理想气体状态方程可知状态 B 的温度高于状态 D 的温度，则状态 B 气体分子热运动的平均动能大，故 D 错误。

故选 C。
}

\item
%% number: 5
%% paperName: 2025年天津高考第 5 题 5 分
%% typeId: 1
%% type: 单选题
%% chapter: 光学
%% point: 光的干涉
%% method: 无
%% score: 5
%% degree: 550
%% duplicateId: 0
%% body:
在进行双缝干涉实验时，用光传感器测量干涉区域不同位置的光照强度，可以方便地展示干涉条纹分布。a、b 两束单色光分别经过同一双缝干涉装置后，其光照强度与位置的关系如图所示，图中光照强度极大值位置对应亮条纹中心，极小值位置对应暗条纹中心，上下两图中相同横坐标代表相同位置，则\xzanswer{D}

\onepicture[w=6.30cm, h=4.77cm]{figs/05.png}

\fourchoices[answer=D]
{a、b 的光子动量大小相同}
{a 的光子能量小于 b 的光子能量}
{在真空中，a 的波长大于 b 的波长}
{通过同一单缝衍射装置，a 的中央亮条纹窄}

%% answer: D
%% memo:
\memoanswer{
C．由图可知，a 光干涉条纹间距小于 b 光干涉条纹间距，根据条纹间距公式 $\Delta x = \frac{L}{d}\lambda$ 可知 a 光的波长小于 b 光的波长，故 C 错误；

A．根据光子动量公式 $p = \frac{h}{\lambda}$ 可知 a 光光子动量大于 b 光光子动量，故 A 错误；

B．根据 $\varepsilon = \frac{hc}{\lambda}$ 可知 a 光光子能量大于 b 光光子能量，故 B 错误；

D．波长越长，通过单缝衍射装置，中央亮条纹越宽，则通过同一单缝衍射装置，a 的中央亮条纹窄，故 D 正确。

故选 D。
}

\item
%% number: 6
%% paperName: 2025年天津高考第 6 题 5 分
%% typeId: 2
%% type: 多选题
%% chapter: 交流电
%% point: 变压器
%% method: 无
%% score: 5
%% degree: 550
%% duplicateId: 0
%% body:
如图所示，理想变压器的原线圈接在电压有效值不变的正弦交流电源上，副线圈接有滑动变阻器和定值电阻，各电表均为理想交流电表。当滑动变阻器的滑动片向右滑动时，则\xzanswer{CD}

\onepicture[w=7.50cm, h=4.13cm]{figs/06.png}

\fourchoices[answer=CD]
{电流表 A$_{1}$ 的示数变小}
{电流表 A$_{2}$ 的示数变小}
{电压表 V 的示数变大}
{变压器的输入功率变大}

%% answer: CD
%% memo:
\memoanswer{
理想变压器的输出的电压由输入电压和电压比决定，输入电压不变，所以输出电压也不会变：

AB．当滑动变阻器的滑动片向右滑动时，滑动变阻器接入电路的电阻变小，根据欧姆定律 $I = \frac{U}{R_{\text{总}}}$ 可知副线圈回路中电流变大，即电流表 A$_{2}$ 的示数变大，根据 $\frac{I_{1}}{I_{2}} = \frac{n_{2}}{n_{1}}$ 可知原线圈回路中电流变大，即电流表 A$_{1}$ 的示数变大，故 AB 错误；

C．电压表 V 测电阻 $R_{0}$ 两端的电压，滑动变阻器接入电路的电阻变小，根据串联电路分压规律可知电压表 V 的示数变大，故 C 正确；

D．根据 $P = \frac{U^{2}}{R}$ 可知副线圈回路的功率变大，根据理想变压器原副线圈功率相等可知变压器的输入功率变大，故 D 正确。

故选 CD。
}

\item
%% number: 7
%% paperName: 2025年天津高考第 7 题 5 分
%% typeId: 2
%% type: 多选题
%% chapter: 机械波
%% point: 波的图象
%% method: 无
%% score: 5
%% degree: 550
%% duplicateId: 0
%% body:
位于坐标原点的波源从平衡位置开始沿 $y$ 轴运动，在均匀介质中形成了一列沿 $x$ 轴正方向传播的简谐波，P 和 Q 是平衡位置分别位于 $x = 3 \Um$ 和 $x = 7 \Um$ 处的两质点，$t = 0$ 时波形如图所示，此时 Q 刚开始振动，$t = 1 \Us$ 时 Q 第一次到达波谷。则\xzanswer{AB}

\onepicture[w=7.00cm, h=3.88cm]{figs/07.png}

\fourchoices[answer=AB]
{该波在此介质中的波速为 $2 \Ums$}
{波源开始振动时的运动方向沿 $y$ 轴负方向}
{P 的位移随时间变化的关系式为 $y = \sin \pi t \Ucm$}
{平衡位置位于 $x = 9 \Um$ 处的质点在 $t = 6 \Us$ 时第一次到达波峰}

%% answer: AB
%% memo:
\memoanswer{
AB．由图根据同侧法可知，Q 点开始振动时的运动方向沿 $y$ 轴负方向，则波源开始振动时的运动方向沿 $y$ 轴负方向，由图可知，该波的波长为 $\lambda = 2 \times (7 - 3) \Um = 8 \Um$

根据题意可知，周期 $T = 4t = 4 \Us$

则该波在此介质中的波速为 $v = \frac{\lambda}{T} = 2 \Ums$，故 AB 正确；

C．根据公式可得 $\omega = \frac{2\pi}{T} = \frac{1}{2}\pi \Urads$

由同侧法可知，此时 P 点沿 $y$ 轴正方向振动，振幅为 $1 \Ucm$，则 P 的位移随时间变化的关系式为 $y = \sin \frac{\pi}{2}t \Ucm$，故 C 错误；

D．该波传播到 $x = 9 \Um$ 处的时间为 $t' = \frac{9 - 7}{2} \Us = 1 \Us$

质点起振方向沿 $y$ 轴负方向，则第一次到达波峰的时间为 $t'' = \frac{3}{4}T = 3 \Us$

则平衡位置位于 $x = 9 \Um$ 处的质点第一次到达波峰的时间为 $t = t' + t'' = 4 \Us$，故 D 错误。

故选 AB。
}

\item
%% number: 8
%% paperName: 2025年天津高考第 8 题 5 分
%% typeId: 2
%% type: 多选题
%% chapter: 电场
%% point: 电场线
%% method: 无
%% score: 5
%% degree: 450
%% duplicateId: 0
%% body:
如图所示，在一固定点电荷形成的电场中，一试探电荷仅在静电力作用下先后经过 a、b 两点，图中箭头方向表示试探电荷在 a、b 两点处的受力方向，则\xzanswer{BC}

\onepicture[w=9.00cm, h=2.55cm]{figs/08.png}

\fourchoices[answer=BC]
{a 点电势一定高于 b 点电势}
{试探电荷与场源电荷电性一定相同}
{a 点电场强度一定大于 b 点电场强度}
{试探电荷的电势能一定先减小后增大}

%% answer: BC
%% memo:
\memoanswer{
AB．根据点电荷电场分布特点和试探电荷受力情况，试探电荷轨迹如图所示。

\onepicture[w=7.50cm, h=3.43cm]{figs/08_an.png}

可知，点电荷位于 O 点，两电荷带同种电荷，由于电性无法判断，所以点电荷周围电场方向无法判断，a 点电势和 b 点电势高低无法判断，故 A 错误，B 正确；

C．由于 a 点离点电荷较近，由 $E = \frac{kQ}{r^{2}}$ 可知，a 点电场强度一定大于 b 点电场强度，故 C 正确；

D．由轨迹图可知，电场力方向与运动方向的夹角先为钝角，后为锐角，所以电场力先做负功，后做正功，试探电荷的电势能一定先增大后减小，故 D 错误；

故选 BC。
}

\item
%% number: 9
%% paperName: 2025年天津高考第 9 题 6 分
%% typeId: 4
%% type: 实验题
%% chapter: 曲线运动
%% point: 研究平抛运动实验
%% method: 无
%% score: 6
%% degree: 650
%% duplicateId: 0
%% body:
某实验小组探究平抛运动的特点。
\begin{enumerate}
\item
实验采用图~\ref{2025天津09a}所示装置，将白纸和复写纸重叠并固定在竖直背板上，钢球在斜槽中某一高度滚下，从末端水平飞出，落在挡板上，在白纸上挤压出一个痕迹点。移动挡板，重复实验，在白纸上留下若干痕迹点。为正确研究钢球运动的规律，将白纸从背板上取下前还应根据\tkanswer{铅垂线}在白纸上标记竖直方向；
\item
利用手机和计算机可以方便地记录钢球做平抛运动的轨迹并分析其运动规律。该实验小组利用视频处理软件分析钢球某次平抛运动的录制视频，得到的水平位移和竖直位移随时间变化的图像如图~\ref{2025天津09b}所示。图中图线\tkanswer{甲}（选填“甲”或“乙”）为水平位移随时间变化的图线，此次钢球做平抛运动的初速度为\tkanswer{2}$\Ums$。
\end{enumerate}

\twopicture[labelA=2025天津09a, labelB=2025天津09b, wA=4.20cm, wB=5.70cm, align=right]{figs/09a.png}{figs/09b.png}

%% answer:
\jdanswer{
\begin{enumerate}
\item
铅垂线
\item
甲，$2$
\end{enumerate}
}
%% memo:
\memoanswer{
（1）为正确研究钢球运动的规律，将白纸从背板上取下前还应根据铅垂线在白纸上标记竖直方向。

（2）水平方向，钢球做匀速直线运动，所以位移随时间均匀增大，即图线甲为水平位移随时间变化的图线，此次钢球做平抛运动的初速度为
$v = \frac{0.3}{0.15} \Ums = 2 \Ums$
}

\item
%% number: 9
%% paperName: 2025年天津高考第 9 题 6 分
%% typeId: 4
%% type: 实验题
%% chapter: 电路
%% point: 伏安法测电阻
%% method: 无
%% score: 6
%% degree: 550
%% duplicateId: 0
%% body:
某同学用伏安法测量电阻 $R_{x}$ 的阻值（约为 $5 \UO$）。除滑动变阻器、电源（电动势为 $3 \UV$，内阻不计）、开关、导线外，还有下列器材可供选择：

A．电压表（量程 $0 \sim 3 \UV$，内阻约为 $3 \UkO$）

B．电压表（量程 $0 \sim 15 \UV$，内阻约为 $15 \UkO$）

C．电流表（量程 $0 \sim 0.6 \UA$，内阻约为 $0.1 \UO$）

D．电流表（量程 $0 \sim 3 \UA$，内阻约为 $0.02 \UO$）

\begin{enumerate}
\item
为使测量尽量准确，实验中电压表应选\tkanswer{A}，电流表应选\tkanswer{C}（均填器材前的字母代号）；
\item
该同学连接好的电路如图~\ref{2025天津09c}所示，闭合开关后，电流表和电压表的读数均为零，经排查，发现电路中的七条导线中混进了一条内部断开的导线。为了确定哪条导线断开，该同学用多用电表的直流电压挡测量电源正极 a 与接线柱 b 间电压，读数为零；测量 a 与接线柱 c 间电压，读数不为零，则可以确定导线\tkanswer{6}内部断开（填图中导线编号）；
\item
更换导线后进行实验，该同学测得多组电压数据 $U$ 和电流数据 $I$，描绘出 $U-I$ 图像如图~\ref{2025天津09d}所示，测得图中直线的斜率为 $k$，已知电压表的内阻为 $R_{\mathrm{V}}$，若考虑电压表内阻的影响，则 $R_{x}$ 阻值的表达式为\tkanswer{$\frac{kR_{\mathrm{V}}}{R_{\mathrm{V}} - k}$}。
\end{enumerate}

\twopicture[labelA=2025天津09c, labelB=2025天津09d, wA=8.00cm, wB=4.80cm, align=right]{figs/09c.png}{figs/09d.png}

%% answer:
\jdanswer{
\begin{enumerate}
\item
A，C
\item
6
\item
$\frac{kR_{\mathrm{V}}}{R_{\mathrm{V}} - k}$
\end{enumerate}
}
%% memo:
\memoanswer{
（1）电源电动势为 $3 \UV$，为了减小误差，电压表选用量程为 $3 \UV$ 的 A；

流过电阻的最大电流约为 $I_{max} = \frac{E}{R_{x}} = \frac{3}{5} \UA = 0.6 \UA$，为了减小误差，电流表选用量程为 $0.6 \UA$ 的 C。

（2）电流表和电压表示数都为零，在不可能是 4、5 导线内部断开，该同学用多用电表的直流电压挡测量电源正极 a 与接线柱 b 间电压，读数为零；测量 a 与接线柱 c 间电压，读数不为零，则可以确定导线 6 内部断开。

（3）根据电路图可得
$I = \frac{U}{R_{\mathrm{V}}} + \frac{U}{R_{x}}$

整理可得
$U = \frac{R_{\mathrm{V}}R_{x}}{R_{x} + R_{\mathrm{V}}}I$

则 $\frac{R_{\mathrm{V}}R_{x}}{R_{x} + R_{\mathrm{V}}} = k$

可得 $R_{x} = \frac{kR_{\mathrm{V}}}{R_{\mathrm{V}} - k}$
}

\item
%% number: 10
%% paperName: 2025年天津高考第 10 题 14 分
%% typeId: 6
%% type: 计算题
%% chapter: 运动和力
%% point: 动量和能量
%% method: 无
%% score: 14
%% degree: 650
%% duplicateId: 0
%% body:
如图所示，半径为 $R = 0.45 \Um$ 的四分之一圆轨道 AB 竖直固定放置，与水平桌面在 B 点平滑连接。质量为 $m = 0.12 \Ukg$ 的玩具小车从 A 点由静止释放，运动到桌面上 C 点时与质量为 $M = 0.18 \Ukg$ 的静置物块发生碰撞并粘在一起，形成的组合体匀减速滑行 $x = 0.20 \Um$ 至 D 点停止。A 点至 C 点光滑，小车和物块碰撞时间极短，小车、物块及组合体均视为质点，$g$ 取 $10 \Umsq$，不计空气阻力。求：
\begin{enumerate}
\item
小车运动至圆轨道 B 点时所受支持力 $F_{N}$ 的大小；
\item
小车与物块碰撞后瞬间组合体速度 $v$ 的大小；
\item
组合体与水平桌面 CD 间的动摩擦因数 $\mu$ 的值。
\end{enumerate}

\onepicture[w=7.00cm, h=3.53cm, align=right]{figs/10.png}

%% answer:
\jdanswer{
\begin{enumerate}
\item
$F_{N} = 3.6 \UN$
\item
$v = 1.2 \Ums$
\item
$\mu = 0.36$
\end{enumerate}
}
%% memo:
\memoanswer{
（1）设小车运动至圆轨道 B 点时的速度大小为 $v_{0}$，由机械能守恒定律，有
$mgR = \frac{1}{2}mv_{0}^{2}$

由牛顿第二定律，有
$F_{N} - mg = m\frac{v_{0}^{2}}{R}$

代入数据，联立解得 $F_{N} = 3.6 \UN$

（2）小车与物块在 C 点碰撞，在水平方向由动量守恒，有
$mv_{0} = (M + m)v$

代入数据，联立解得 $v = 1.2 \Ums$

（3）组合体水平方向受动摩擦力作用，从 C 点匀减速运动至 D 点静止，由动能定理，有
$-\mu(M + m)gx = 0 - \frac{1}{2}(M + m)v^{2}$

代入数据，联立解得 $\mu = 0.36$
}

\item
%% number: 11
%% paperName: 2025年天津高考第 11 题 16 分
%% typeId: 6
%% type: 计算题
%% chapter: 磁场
%% point: 带电粒子在组合场中的运动
%% method: 无
%% score: 16
%% degree: 550
%% duplicateId: 0
%% body:
如图所示，纸面内水平虚线下方存在竖直向上的匀强电场，虚线上方存在垂直于纸面的匀强磁场。一质量为 $m$、电荷量为 $q$ 的粒子从电场中的 O 点以水平向右的速度开始运动，在静电力的作用下从 P 点进入磁场，射入磁场时的速度大小为 $v$、方向与竖直方向夹角为 $\theta$，粒子返回电场前的运动轨迹过 P 点正上方的 Q 点，P、Q 间距离及 O、P 间的水平距离均为 $L$。不计粒子重力。
\begin{enumerate}
\item
判断粒子的电性；
\item
求电场强度大小 $E$；
\item
求磁感应强度大小 $B$。
\end{enumerate}

\onepicture[h=5.60cm, align=right]{figs/11.png}

%% answer:
\jdanswer{
\begin{enumerate}
\item
正电
\item
$E = \frac{mv^{2}\sin\theta\cos\theta}{qL}$
\item
$B = \frac{2mv\sin\theta}{qL}$
\end{enumerate}
}
%% memo:
\memoanswer{
（1）根据题意可知，粒子向上偏转，所受电场力向上，与电场方向相同，则粒子带正电。

（2）设粒子在电场中运动的时间为 $t$，水平方向上由运动学公式，有
$L = vt\sin\theta$

设粒子在电场中运动的加速度为 $a$，由牛顿第二定律，有
$qE = ma$

竖直方向上由运动学公式，有
$v\cos\theta = at$

联立上述各式，得
$E = \frac{mv^{2}\sin\theta\cos\theta}{qL}$

（3）设粒子在磁场中做圆周运动的轨迹半径为 $r$，由几何关系，得
$L = 2r\sin\theta$

洛伦兹力提供向心力，有
$qvB = m\frac{v^{2}}{r}$

联立得
$B = \frac{2mv\sin\theta}{qL}$
}

\item
%% number: 12
%% paperName: 2025年天津高考第 12 题 18 分
%% typeId: 6
%% type: 计算题
%% chapter: 电磁感应
%% point: 电磁感应受力分析
%% method: 微元
%% score: 18
%% degree: 450
%% duplicateId: 0
%% body:
轴向磁通风力发电机在新能源领域中有广泛应用，其原理可简化为一圆盘发电机。如图所示，发电机的中心轴为固定不动的圆柱，一外半径为 $3l$、厚度均匀的环形导体盘套在轴上，接触良好并可绕轴转动，导体盘轴线与中心轴的轴线重合。整个装置处在方向与轴线平行的匀强磁场中，磁感应强度大小为 $B$。在风力的作用下，导体盘以角速度 $\omega$ 匀速转动，导体盘的内、外缘为发电机的两个电极。两极接在外电阻两端后，导体盘上各处均有沿半径流动的电流。
\begin{enumerate}
\item
磁场方向与导体盘转动方向如图所示，试判断导体盘的外缘是发电机的正极还是负极；
\item
若外电阻阻值为 $R$，导体盘电阻忽略不计、内半径为 $l$，通过导体盘上相同圆心角区域内的电流相同。求作用在导体盘上圆心角为 $\theta$ 区域（$\theta$ 很小，可视为导体棒）上的安培力大小 $F$ 与 $\theta$ 的关系式；
\item
若外电阻阻值忽略不计，导体盘电阻不可忽略，距离轴线为 $r$ 处的电阻率 $\rho$ 与 $r$ 成正比，比例系数为 $k$，即 $\rho = kr$，导体盘厚度为 $d$、内半径大小可调。求导体盘发热功率最大时内半径的大小。
\end{enumerate}

\onepicture[w=8.50cm, h=5.72cm, align=right]{figs/12.png}

%% answer:
\jdanswer{
\begin{enumerate}
\item
正极
\item
$F = \frac{4B^{2}l^{3}\omega}{\pi R}\theta$
\item
$x = l$
\end{enumerate}
}
%% memo:
\memoanswer{
（1）根据右手定则可知导体盘的外缘是发电机的正极。

（2）设导体盘上圆心角为 $\theta$ 的区域切割磁场的平均速度为 $v$，有
$v = \frac{1}{2}(3l + l)\omega$

设导体盘的电动势为 $E$，由法拉第电磁感应定律，有
$E = B \cdot 2lv$

设回路中的总电流为 $I$，由闭合电路的欧姆定律，有
$I = \frac{E}{R}$

设导体盘上圆心角为 $\theta$ 区域的电流为 $I'$，则
$I' = \frac{\theta}{2\pi}I$

作用在导体盘上圆心角为 $\theta$ 区域上的安培力大小为
$F = BI' \cdot 2l$

联立上述各式，得
$F = \frac{4B^{2}l^{3}\omega}{\pi R}\theta$

（3）如图所示，设距离轴线为 $r$ 处的沿半径方向的微小长度为 $\Delta r$、横截面积为 $\theta rd$ 的导体电阻为 $\Delta R'$，有

\onepicture[w=7.00cm, h=2.23cm]{figs/12_an.png}

$\Delta R' = kr\frac{\Delta r}{\theta rd}$

设导体盘的内半径大小为 $x$、圆心角为 $\theta$ 区域的电阻为 $R'$，有
$R' = k\frac{3l - x}{\theta d}$

设导体盘的电阻为 $R_{\text{盘}}$，则
$R_{\text{盘}} = \frac{\theta R'}{2\pi}$

设内半径为 $x$ 时导体盘的电动势为 $E'$，可知
$E' = B(3l - x)\left[\frac{1}{2}(3l + x)\omega\right]$

设导体盘发热功率为 $P$，有
$P = \frac{E'^{2}}{R_{\text{盘}}}$

联立上述式子可得
$P = \frac{\pi dB^{2}\omega^{2}(3l - x)(3l + x)^{2}}{2k}$

根据数学求导部分知识可得当 $x = l$ 时，导体盘发热功率 $P$ 最大。
}

\end{enumerate}

\end{document}
